%------------------------------------------------------------------------------%

\usepackage{integracao_e_series}

\title{Integração por Substituição Simples}

%------------------------------------------------------------------------------%
\begin{document}

\addtitle

%------------------------------------------------------------------------------%
\section{Integração por Substituição Simples}

%------------------------------------------------------------------------------%
\begin{frame}{Justificativa}
  \onslide<+->{
    Assumindo que \(F'(x)=f(x)\) e \(g'(x)\) é a derivada e \(g(x)\)
  }

  \onslide<+->{temos}
  \begin{align*}
    \onslide<2->{\Big(F\big(g(x)\big)\Big)'
    & \;=\; f\big(g(x)\big) g'(x) \\[1mm]}
    \onslide<+->{{\color<4>{magenta}\int \Big(}F\big(g(x)\big){\color<4>{magenta}\Big)' dx}
    & \;=\; \int f\big(g(x)\big) g'(x) \, dx \\[1mm]}
    \onslide<+->{F\big(g(x)\big)
    & \;=\; \int f\big(g(x)\big) g'(x) \, dx \\}
    \onslide<+->{\int f\big(g(x)\big) g'(x) \, dx
    & \;=\; F\big(g(x)\big)}
    \onslide<+->{\;=\; \int f({\color<6>{blue}u})\,{\color<6>{blue}du}}
  \end{align*}
  \onslide<6->{onde fizemos \({\color<6>{blue}u} = g(x)\) e \({\color<6>{blue}du} = g'(x)\,dx\)}
\end{frame}

%------------------------------------------------------------------------------%
\begin{frame}{Integração por Substituição Simples}
  Regra
  \[
    \int f\big(g(x)\big) g'(x) \, dx
    = \int f(u)\,{du}
  \]
  onde \(u = g(x)\) e \(du = g'(x)\,dx\)

  \pause
  \vspace{\baselineskip}
  Integrais definidas
  \[
    \int_a^b f\big(g(x)\big) g'(x) \, dx
    = \int_{u(a)}^{u(b)} f(u)\,{du}
  \]
\end{frame}

%------------------------------------------------------------------------------%
\begin{frame}{Casos Particulares}
  Se\;
  \(
    \int f(x) dx = F(x) + C
  \)\; e \(k\) é uma constante

  \pause
  \vspace{\baselineskip}
  \[ \int f(kx)  dx \;=\; \frac{1}{\,k\,}F(kx)   + C \]

  \pause
  \[ \int f(kx+b)dx \;=\; \frac{1}{\,k\,}F(kx+b) + C \]
\end{frame}

% 1
%------------------------------------------------------------------------------%
\addtocounter{exemplo}{1}
\section{Exemplo \theexemplo}

%------------------------------------------------------------------------------%
\begin{frame}{Exemplo \theexemplo}
  \onslide<+->{Encontre \(\int 4x^3 \cos\left(x^4+2\right)dx\)}
  \onslide<+->{\( = \int  {\color<3>{blue}\cos\left({\color<4,5>{magenta}x^4+2}\right)}
    {\color<5>{teal}4x^3}\, dx  \)}

  \vspace{\baselineskip}
  \onslide<+->{Identificando as funções}
  \begin{align*}
    \onslide<3->{f  & = {\color<3>{blue}\cos\left({\color<4,5>{magenta}x^4+2}\right)} \\[2mm]}
    \onslide<+->{   & = \cos\big({\color<4,5>{magenta}g(x)}\big)     \\[2mm]}
    \onslide<+->{g  & = {\color<5>{magenta}x^4+2}                  \\[2mm]}
    \onslide<5->{g' & = {\color<5>{teal}4x^3}}
  \end{align*}
\end{frame}

%------------------------------------------------------------------------------%
\begin{frame}{Exemplo \theexemplo}
  \onslide<+->{Nova variável
  \[
    {\color<2>{magenta}u  = x^4+2}
    \qquad\text{e}\qquad
    {\color<2>{blue}du = 4x^3\,dx}
  \]}
  \vspace{-\baselineskip}
  \begin{align*}
    \onslide<+->{\int \cos\left({\color<2>{magenta}x^4+2}\right) \,{\color<2>{blue}4x^3dx}
      & = \int \cos({\color<2>{magenta}u}) {\color<2>{blue}du}                       \\}
    \onslide<+->{
      & = \sin(u) + C                          \\[2mm]}
    \onslide<+->{
      & = \sin\left( {\color{magenta}x^4 + 2} \right) + C}
  \end{align*}
\end{frame}

% 2
%------------------------------------------------------------------------------%
\addtocounter{exemplo}{1}
\section{Exemplo \theexemplo}

%------------------------------------------------------------------------------%
\begin{frame}{Exemplo \theexemplo}
  \onslide<+->{Calcule \(\int \frac{x}{\sqrt{1-4x^2}} dx\)}
  \onslide<+->{\(= \int \left(1-4x^2\right)^{-\sfrac{1}{2}} \,x\,dx\)}

  \vspace{\baselineskip}
  \onslide<+->{Identificando as funções}
  \begin{align*}
    \onslide<3->{f  & = \frac{1}{\sqrt{1-4x^2}} = \left(1-4x^2\right)^{-\sfrac{1}{2}}\\[2mm]}
    \onslide<+->{   & = \big(g(x)\big)^{-\sfrac{1}{2}}\\[2mm]}
    \onslide<+->{g  & = 1-4x^2\\[2mm]
                 g' & = -8x}
  \end{align*}
\end{frame}

%------------------------------------------------------------------------------%
\begin{frame}{Exemplo \theexemplo}
  \onslide<+->{Substituição
  \[
    u = 1-4x^2
    \qquad\text{e}\qquad
    du  =-8x dx
    \quad\Leftrightarrow\quad
    xdx = -\frac{1}{8} du
  \]}
  \begin{align*}
    \onslide<+->{\int \frac{x}{\sqrt{1-4x^2}}dx & = -\frac{1}{8} \int \frac{1}{\sqrt{u}}du \\[2mm]}
    \onslide<+->{                               & = -\frac{1}{8} \int u^{-\sfrac{1}{2}} du \\[2mm]}
    \onslide<+->{                               & = -\frac{1}{8} \left(2 \sqrt{u}\right)+C \\[2mm]}
    \onslide<+->{                               & = -\frac{1}{4}\sqrt{1-4x^2}+C}
  \end{align*}
\end{frame}

% 3
%------------------------------------------------------------------------------%
\addtocounter{exemplo}{1}
\section{Exemplo \theexemplo}

%------------------------------------------------------------------------------%
\begin{frame}{Exemplo \theexemplo}
  \onslide<+->{Calcule \(\int \tan(x) dx\)}
  \onslide<+->{\(= \int \frac{\sin(x)}{\cos(x)} dx\)}
  \onslide<+->{\(= \int \big(\cos(x)\big)^{-1} \sin(x)\,dx\)}

  \vspace{\baselineskip}
  \onslide<+->{Identificando as funções}
  \begin{align*}
    \onslide<4->{f  & = \big(\cos(x)\big)^{-1}\\[2mm]}
    \onslide<+->{   & = \big(g(x)\big)^{-1}\\[2mm]}
    \onslide<+->{g  & = \cos(x) \\[2mm]
                 g' & = -\sin(x)}
  \end{align*}
\end{frame}

%------------------------------------------------------------------------------%
\begin{frame}{Exemplo \theexemplo}
  \onslide<+->{Substituição
  \[
    u = \cos(x)
    \qquad\text{e}\qquad
    du = -\sin(x)dx
    \quad\Leftrightarrow\quad
    -du = \sin(x)dx
  \]}
  \begin{align*}
    \onslide<+->{\int \tan(x) dx & = \int \big(\cos(x)\big)^{-1} \sin(x)\,dx \\}
    \onslide<+->{                & = \int u^{-1} (-du)                       \\}
    \onslide<+->{                & = - \int u^{-1} du                        \\}
    \onslide<+->{                & = - \ln \abs{u} + C                       \\[2mm]}
    \onslide<+->{                & = - \ln \abs{\cos(x)} + C}
  \end{align*}
\end{frame}

% 4
%------------------------------------------------------------------------------%
\addtocounter{exemplo}{1}
\section{Exemplo \theexemplo}

%------------------------------------------------------------------------------%
\begin{frame}{Exemplo \theexemplo}
  \onslide<+->{Encontre o valor médio de \;
  \(
    f(x)=\frac{\ln(x)}{x}
  \)
  \; no intervalo $[1, e]$}

  \vspace{2\baselineskip}
  \onslide<+->{Valor médio
  \[
    f_{\text{médio}}
    = \frac{1}{b-a} \int_a^b f(x)\, dx
  \]}
\end{frame}

%------------------------------------------------------------------------------%
\begin{frame}{Exemplo \theexemplo\ -- Integral Indefinida}
  \onslide<+->{Calcular \(\int \frac{\ln(x)}{x} dx\)}
  \onslide<+->{\(= \int \ln(x) x^{-1} dx\)}

  \vspace{0.5\baselineskip}
  \onslide<+->{Substituição \quad
  \(
    u = \ln(x)
    \qquad\text{e}\qquad
    du = x^{-1} dx
  \)}

  \vspace{0.5\baselineskip}
  \begin{align*}
    \onslide<+->{\int \frac{\ln(x)}{x} dx & = \int u\, du            \\[2mm]}
    \onslide<+->{                         & = \frac{u^2}{2}+C        \\[2mm]}
    \onslide<+->{                         & = \frac{1}{2}\big(\ln(x)\big)^2 + C}
  \end{align*}
\end{frame}

%------------------------------------------------------------------------------%
\begin{frame}{Exemplo \theexemplo\ -- Integral Definida}
  \begin{align*}
    \onslide<+->{\int_1^e \frac{\ln(x)}{x} dx
      & = \frac{1}{2}\big(\ln(x)\big)^2 \evalat{1}{e}                 \\[2mm]}
    \onslide<+->{
      & = \frac{\big(\ln(e)\big)^2}{2} - \frac{\big(\ln(1)\big)^2}{2} \\[2mm]}
    \onslide<+->{
      & = \frac{1}{2}}
  \end{align*}
\end{frame}

%------------------------------------------------------------------------------%
\begin{frame}{Exemplo \theexemplo\ -- Valor Médio}
  \begin{align*}
    \onslide<+->{f_{\text{médio}} & = \frac{1}{b-a} \int_a^b f(x)\, dx           \\[2mm]}
    \onslide<+->{                 & = \frac{1}{e-1} \int_1^e \frac{\ln(x)}{x} dx \\[2mm]}
    \onslide<+->{                 & = \frac{1}{e-1} \frac{1}{2}                  \\[2mm]}
    \onslide<+->{                 & = \frac{1}{2(e-1)}}
  \end{align*}
\end{frame}

% 6
%------------------------------------------------------------------------------%
\addtocounter{exemplo}{1}
\section{Exemplo \theexemplo}

%------------------------------------------------------------------------------%
\begin{frame}{Exemplo \theexemplo}
  Determine se a integral \(\int_{-\infty}^{\infty} x^3e^{-x^4} dx\)
  é convergente ou divergente

  \pause
  \vspace{\baselineskip}
  Integral Imprópria
  \[
    \int_{-\infty}^{\infty} x^3e^{-x^4} dx
    \;=\; \lim_{r \to-\infty} \int_{r}^{0} x^3e^{-x^4} dx
    \;+\; \lim_{s \to \infty} \int_{0}^{s} x^3e^{-x^4} dx
  \]
\end{frame}

%------------------------------------------------------------------------------%
\begin{frame}{Exemplo \theexemplo\ -- Integral Indefinida}
  \onslide<+->{Calcular \(\int x^3e^{-x^4} dx\)}

  \onslide<+->{\vspace{0.1\baselineskip}
  Substituição \quad
  \(
    u  = -x^4
    \qquad\text{e}\qquad
    du = -4x^3dx
    \quad\Leftrightarrow\quad
    x^3dx = \frac{du}{-4}
  \)}

  \vspace{0.1\baselineskip}
  \begin{align*}
    \onslide<+->{\int x^3e^{-x^4} dx & = \int e^u \left(-\frac{du}{4}\right) \\[1mm]}
    \onslide<+->{                    & = -\frac{1}{4}\int e^u du             \\[1mm]}
    \onslide<+->{                    & = -\frac{1}{4}e^u+C                   \\[1mm]}
    \onslide<+->{                    & = -\frac{1}{4}e^{-x^4}+C}
  \end{align*}
\end{frame}

%------------------------------------------------------------------------------%
\begin{frame}{Exemplo \theexemplo\ -- Integral Imprópria 1}
\begin{align*}
  \onslide<+->{\int_{-\infty}^{0} x^3e^{-x^4} dx
                 & = \lim_{r \to -\infty} \int_r^0 x^3e^{-x^4} dx                \\[3mm]}
  \onslide<+->{  & = \lim_{s \to -\infty} \left(-\frac{e^{-x^4}}{4}\right) \evalat{r}{0} \\[3mm]}
  \onslide<+->{  & = \lim_{r \to -\infty} \left(-\frac{1}{4}\left( e^0 - e^{-r^4} \right)\right) \\[3mm]}
  \onslide<+->{  & = -\frac{1}{4}\left( 1 - \lim_{r \to -\infty} e^{-r^4} \right)               \\[3mm]}
  \onslide<+->{  & = -\frac{1}{4}}
\end{align*}
\end{frame}

%------------------------------------------------------------------------------%
\begin{frame}{Exemplo \theexemplo\ -- Integral Imprópria 2}
  \begin{align*}
    \onslide<+->{\int_{0}^{\infty} x^3e^{-x^4} dx
                 & = \lim_{s \to \infty} \int_0^s x^3e^{-x^4} dx                           \\[3mm]}
    \onslide<+->{& = \lim_{s \to \infty} \left(-\frac{e^{-x^4}}{4}\right) \evalat{0}{s} \\[3mm]}
    \onslide<+->{& = \lim_{s \to \infty} \left(-\frac{1}{4}\left( e^{-s^4} - e^0 \right)\right) \\[3mm]}
    \onslide<+->{& = -\frac{1}{4}\left( \lim_{s \to \infty}e^{-s^4} - 1 \right)                \\[3mm]}
    \onslide<+->{& = \frac{1}{4}}
  \end{align*}
\end{frame}

%------------------------------------------------------------------------------%
\begin{frame}{Exemplo \theexemplo\ -- Conclusão}
  \begin{align*}
  \onslide<+->{\int_{-\infty}^{\infty} x^3e^{-x^4} dx
   & = \lim_{r \to-\infty} \int_{r}^{0} x^3e^{-x^4} dx
     + \lim_{s \to \infty} \int_{0}^{s} x^3e^{-x^4} dx \\[3mm]}
  \onslide<+->{& = -\frac{1}{4} +\frac{1}{4}  \\[3mm]}
  \onslide<+->{& = 0}
\end{align*}
\onslide<+->{A integral converge para zero}
\end{frame}

%------------------------------------------------------------------------------%
\section{Lista Mínima}

%------------------------------------------------------------------------------%
\begin{frame}{Lista Mínima}
  Estudar a Seção 4.2 da Apostila

  Exercícios: 1a-f, 2, 4, 5a-f

  \vfill
  Atenção: A prova é baseada no livro, não nas apresentações
\end{frame}



%------------------------------------------------------------------------------%
\end{document}
%------------------------------------------------------------------------------%
